\section*{अध्याय \ref{ch:g6:inside-the-earth} --- पृथ्वी के अंदर}

\begin{solution}{exo:g6:inside-the-earth:1}
\omterm{def:g6:inside-the-earth:layers}{भूपर्पटी}, \omterm{def:g6:inside-the-earth:layers}{प्रावार}, \omterm{def:g6:inside-the-earth:layers}{बाह्य क्रोड}, \omterm{def:g6:inside-the-earth:layers}{आंतरिक क्रोड}। \omterm{def:g6:inside-the-earth:layers}{बाह्य क्रोड}
द्रव है। बाह्य और \omterm{def:g6:inside-the-earth:layers}{आंतरिक क्रोड} धातु के बने हैं (ज़्यादातर लोहा, थोड़े
निकल के साथ); \omterm{def:g6:inside-the-earth:layers}{भूपर्पटी} और \omterm{def:g6:inside-the-earth:layers}{प्रावार} \omterm{def:g3:stones-pebbles-sand:rock}{चट्टान} के बने हैं।
\end{solution}

\begin{solution}{exo:g6:inside-the-earth:2}
(क) ग़लत: \omterm{def:g6:inside-the-earth:layers}{भूपर्पटी} \omterm{def:g3:continents-and-oceans:continent}{महाद्वीपों} के नीचे (लगभग 30~km) \omterm{def:g2:water-on-earth:ocean}{महासागरों} के नीचे (लगभग 5~km)
से ज़्यादा मोटी है। (ख) ग़लत: \omterm{def:g6:inside-the-earth:layers}{क्रोड} धातु का बना है, लोहे और निकल का एक
मिश्रधातु। (ग) सही। (घ) सही: 12\,262~m, लगभग 12.3~km।
\end{solution}

\begin{solution}{exo:g6:inside-the-earth:3}
(क) 2891~km पर, लगभग 2900~km। (ख) 5150~km पर। (ग) $5150 - 2891 = 2259$:
लगभग 2260~km।
\end{solution}

\begin{solution}{exo:g6:inside-the-earth:4}
``छिलका \underline{\omterm{def:g6:inside-the-earth:layers}{भूपर्पटी}} जैसा है, सफ़ेदी
\underline{\omterm{def:g6:inside-the-earth:layers}{प्रावार}} जैसी है और ज़र्दी \underline{\omterm{def:g6:inside-the-earth:layers}{क्रोड}} जैसी है। पर
ज़र्दी के उलट, \underline{\omterm{def:g6:inside-the-earth:layers}{क्रोड}} के दो हिस्से हैं: एक द्रव \underline{\omterm{def:g6:inside-the-earth:layers}{बाह्य क्रोड}}
और एक ठोस \underline{\omterm{def:g6:inside-the-earth:layers}{आंतरिक क्रोड}}।''
\end{solution}

\begin{solution}{exo:g6:inside-the-earth:5}
$8.849 \div 6371 \approx 0.00139$, और $0.00139 \times 100 \approx 0.14$:
एवरेस्ट त्रिज्या का केवल लगभग \textbf{0.14~\%} है। दूर से देखने पर
पृथ्वी एक बहुत चिकना गोला है: उसका सबसे ऊँचा \omterm{def:g1:landscapes-around-us:landform}{पर्वत} एक नन्हा उभार है।
\end{solution}

\begin{solution}{exo:g6:inside-the-earth:6}
(क) $30 \div 5 = 6$: छह गुना मोटी। (ख) $100 \div 5 = 20$: बीस गुना
मोटी।
\end{solution}

\begin{solution}{exo:g6:inside-the-earth:7}
(क) $12 + 1 \times 25 = 37~^\circ$C। (ख) $12 + 1.5 \times 25 = 12 + 37.5 =
49.5~^\circ$C। (ग) $12 + 2 \times 25 = 62~^\circ$C।
\end{solution}

\begin{solution}{exo:g6:inside-the-earth:8}
बढ़त 2~km में $74 - 14 = 60$ डिग्री है, इसलिए $60 \div 2 =
\textbf{30}$~डिग्री प्रति किलोमीटर। हाँ: 30 पूर्वी औसत
(25) और पश्चिमी औसत (34) के बीच है।
\end{solution}

\begin{solution}{exo:g6:inside-the-earth:9}
(क) \omterm{def:g6:inside-the-earth:layers}{क्रोड} के ठीक ऊपर लगभग 13.7~km प्रति सेकंड; उसके ऊपरी सिरे के ठीक नीचे
लगभग 8.1~km प्रति सेकंड: P तरंगें अचानक धीमी हो जाती हैं। (ख)
2891~km और 5150~km के बीच, \omterm{def:g6:inside-the-earth:layers}{बाह्य क्रोड} में। (ग) S तरंगें किसी द्रव में से नहीं
गुज़र सकतीं: \omterm{def:g6:inside-the-earth:layers}{बाह्य क्रोड} द्रव है।
\end{solution}

\begin{solution}{exo:g6:inside-the-earth:10}
(क) $12.3 \div 30 = 0.41$: एक औसत महाद्वीपीय \omterm{def:g6:inside-the-earth:layers}{भूपर्पटी} का लगभग
\textbf{41~\%}। (ख) नहीं: किसी \omterm{def:g3:continents-and-oceans:continent}{महाद्वीप} के नीचे \omterm{def:g6:inside-the-earth:layers}{प्रावार} \omterm{def:g6:inside-the-earth:layers}{भूपर्पटी} के नीचे, औसतन
लगभग 30~km नीचे, शुरू होता है; वेध-छिद्र उससे बहुत ऊपर, \omterm{def:g6:inside-the-earth:layers}{भूपर्पटी} में ही रुक गया।
\end{solution}

\begin{solution}{exo:g6:inside-the-earth:11}
(क) $29 \times 6371 = 184\,759$, और $15 + 184\,759 = 184\,774$: लगभग
185\,000~डिग्री। (ख) प्रयोगशाला \omterm{def:g6:inside-the-earth:layers}{आंतरिक क्रोड} के किनारे पर लगभग 6000~डिग्री
बताती है: लियो की संख्या लगभग 30~गुना ज़्यादा बड़ी है
($184\,774 \div 6000 \approx 31$)। (ग) 29~डिग्री प्रति किलोमीटर की दर
केवल पहले लगभग 2~km में नापी गई थी। लियो उसे मापों से हज़ारों
किलोमीटर आगे काम में लेता है; बहुत नीचे तापमान बहुत ज़्यादा
धीरे बढ़ता है। किसी दर का प्रयोग केवल वहीं किया जा सकता है जहाँ उसे नापा गया हो।
\end{solution}

\begin{solution}{exo:g6:inside-the-earth:12}
(क) अगर \omterm{def:g6:inside-the-earth:layers}{भूपर्पटी} औसत से हल्की है, तो और गहरे का पदार्थ 5.5 से ज़्यादा घना होना
चाहिए, ताकि औसत ऊपर आए। (ख) \omterm{def:g6:inside-the-earth:layers}{बाह्य क्रोड} और
\omterm{def:g6:inside-the-earth:layers}{आंतरिक क्रोड} में (लगभग 9.9 से लगभग 13~ग्राम प्रति घन सेंटीमीटर); \omterm{def:g6:inside-the-earth:layers}{प्रावार}
अपनी एकदम तली पर बस 5.6 तक पहुँचता है। (ग) ज़्यादा घने पदार्थ हल्के पदार्थों के नीचे
डूबते हैं, जैसे पत्थर पानी में डूबता है: भारी धातु बीच में
पहुँच गई और हल्की \omterm{def:g3:stones-pebbles-sand:rock}{चट्टानें} ऊपर।
\end{solution}

\begin{solution}{pb:g6:inside-the-earth:1}
\textbf{1.} $2 \times 6371 = 12\,742$~km। मॉडल 2~m चौड़ा है, यानी
2000~mm।

\textbf{2.} $12\,742 \div 2000 \approx 6.37$: मॉडल का 1~mm असली पृथ्वी के लगभग
6.37~km के बराबर है।

\textbf{3.} $8.849 \div 6.371 \approx 1.39$: मॉडल पर एवरेस्ट लगभग
\textbf{1.4~mm} ऊँचा है, \omterm{def:g3:stones-pebbles-sand:sand}{रेत} का एक कण।

\textbf{4.} $12.3 \div 6.371 \approx 1.9$~mm। मॉडल का केंद्र सतह से
1000~mm (1~m) नीचे है: वेध-छिद्र लगभग 520~गुना छोटा है
($1000 \div 1.9 \approx 520$)।

\textbf{5.} $2891 \div 6.371 \approx 453.8$: \omterm{def:g6:inside-the-earth:layers}{बाह्य क्रोड} का ऊपरी सिरा सतह से
लगभग \textbf{454~mm} नीचे रँगा जाता है।

\textbf{6.} $5150 \div 6.371 \approx 808.3$: सतह से लगभग 808~mm
नीचे। मॉडल पर \omterm{def:g6:inside-the-earth:layers}{आंतरिक क्रोड} की त्रिज्या $1000 - 808 = 192$~mm है।

\textbf{7.} \omterm{def:g6:inside-the-earth:layers}{प्रावार}: लगभग 5~mm से 454~mm तक, लगभग 449~mm। \omterm{def:g6:inside-the-earth:layers}{बाह्य
क्रोड} $808 - 454 = 354$~mm घेरता है और \omterm{def:g6:inside-the-earth:layers}{आंतरिक क्रोड} 192~mm।

\textbf{8.} $30 \div 6.371 \approx 4.7$~mm: एक पतली पट्टी, जिसे एक बारीक
ब्रश से रँगा जा सकता है।

\textbf{9.} $15 + 2 \times 29 = 73~^\circ$C। मॉडल पर 2~km सतह से
$2 \div 6.371 \approx 0.31$~mm नीचे है: मुश्किल से काग़ज़ की
तीन शीटों जितनी मोटाई।

\textbf{10.} $4 \div 6.371 \approx 0.63$~mm, वेध-छिद्र के लगभग 1.9~mm के
मुक़ाबले: खदान उथली है, लगभग $1.9 - 0.6 = 1.3$~mm।

\textbf{11.} कोई पदार्थ पिघलता है या नहीं, यह तापमान और दाब पर निर्भर करता है।
\omterm{def:g6:inside-the-earth:layers}{आंतरिक क्रोड} को सबसे ज़्यादा दबाया जाता है, हवा के दाब के लगभग
33~लाख गुना से: दबा हुआ लोहा केवल ज़्यादा ऊँचे तापमान पर पिघलता है, इसलिए
\omterm{def:g6:inside-the-earth:layers}{आंतरिक क्रोड} का लोहा ठोस रहता है, भले ही वह
\omterm{def:g6:inside-the-earth:layers}{बाह्य क्रोड} से गरम है।

\textbf{12.} $5 \div 6.371 \approx 0.785$: 2-मीटर वाले मॉडल पर महासागरीय \omterm{def:g6:inside-the-earth:layers}{भूपर्पटी}
केवल \textbf{0.78~mm} मोटी रँगी जानी चाहिए, लगभग आठ काग़ज़ की शीटों
जितनी मोटी ($0.78 \div 0.1 \approx 8$)।
\end{solution}
